\section{Problem and Contributions}
\label{sec:overview}
\subsection{Task and motivation}
\citet{xu2008csa} study how to compress tensor-valued objects while retaining
information for reconstruction and recognition. Vectorized image features can
be high-dimensional relative to the training set: their ORL recognition setting
uses $28\times23\times40=25{,}760$ Gabor values per image but only two training
images per identity, or 80 in total. The challenge is to learn a useful compact
representation from limited data, not simply to reshape the input.

\subsection{Central idea and contributions}
CSA learns shared bases along each object's modes and a small core for each
sample. Its contributions are a reconstruction criterion for arbitrary-order
tensors, an alternating eigendecomposition algorithm, a constrained-PCA
interpretation, and reconstruction/recognition experiments on faces and digits.
Unlike Tensorface's organization of an image ensemble by identity, pose, and
illumination, CSA's axes describe an individual object, such as height, width,
and filter channel \citep[Sec.~III-A]{xu2008csa}. This distinction makes axis
definition a modeling choice, rather than treating all tensor representations
as interchangeable.
